% Direct-authority diplomatic French for printed p.128.
\SourcePage{128}

\LayerTarget{tohoku:section:1.4:p4}
\noindent la surjection canonique $B\longrightarrow B/B'$. On définit
dualement l'image directe $u(A')$ d'un quotient $A'$ de $A$, c'est un quotient
de $B$. On démontre pour ces notions toutes les propriétés formelles usuelles.

\LayerTarget{tohoku:section:1.4:exact-sequence}
Enfin, rappelons qu'un couple

\[
  A\mathrel{\mathop{\longrightarrow}^{u}}B
   \mathrel{\mathop{\longrightarrow}^{v}}C
\]

\noindent de deux morphismes consécutifs est dit \textit{exact} si
$\mathrm{Ker}\,v=\mathrm{Im}\,u$, d'où plus généralement la notion de
\textit{suite exacte} de morphismes. Pour qu'une suite

\[
  0\longrightarrow A\longrightarrow B\longrightarrow C
\]

\noindent soit exacte, il faut et il suffit que pour tout $X\in C$, la suite
d'homomorphismes de groupes abéliens suivante soit exacte :

\[
  0\longrightarrow\mathrm{Hom}(X,A)\longrightarrow\mathrm{Hom}(X,B)
  \longrightarrow\mathrm{Hom}(X,C).
\]

\noindent Critère dual pour qu'une suite
$C\longrightarrow B\longrightarrow A\longrightarrow0$ soit exacte. Pour
qu'une suite

\[
  0\longrightarrow A'\longrightarrow A\longrightarrow A''\longrightarrow0
\]

\noindent soit exacte, il faut et il suffit que $u$ soit un monomorphisme et
que $v$ en soit un conoyau généralisé.

\LayerTarget{tohoku:section:1.4:exact-functor}
Soit $F$ un foncteur covariant d'une catégorie abélienne $C$ dans une autre
$C'$. Suivant la terminologie introduite dans [6], nous dirons que $F$ est un
\textit{foncteur semi-exact} (resp. \textit{exact à gauche} resp.
\textit{exact à droite}) si pour toute suite exacte

\[
  0\longrightarrow A'\longrightarrow A\longrightarrow A''\longrightarrow0
\]

\noindent dans $C$, la suite de morphismes correspondants

\[
  0\longrightarrow F(A')\longrightarrow F(A)\longrightarrow F(A'')
  \longrightarrow0
\]

\noindent est exacte en $F(A)$ (resp. en $F(A)$ et $F(A')$, resp. en $F(A)$
et $F(A'')$). $F$ est dit un \textit{foncteur exact} si $F$ est exact à gauche
et à droite, i. e. transforme une suite exacte du type précédent en une suite
exacte ; alors $F$ transforme \textit{toute} suite exacte en une suite exacte.
Si $F$ est exact à gauche, $F$ transforme une suite exacte

\[
  0\longrightarrow A\longrightarrow B\longrightarrow C
\]

\noindent en une suite exacte

\[
  0\longrightarrow F(A)\longrightarrow F(B)\longrightarrow F(C);
\]

\noindent énoncé dual pour les foncteurs exacts à droite. Si $F$ est un
foncteur contravariant, on dit que $F$ est semi-exact (resp. exact à gauche
etc.) si $F$ l'est en tant que foncteur covariant de $C^\circ$ dans $C'$. Le
composé de foncteurs covariants exacts à gauche (resp. à droite) est du même
type. Nous renvoyons à [6] pour d'autres sorites de ce genre, et l'étude des
propriétés d'exactitude des multifoncteurs. Comme exemple important, notons que
$\mathrm{Hom}(A,B)$ est un bifoncteur additif sur $C\times C$, à valeurs dans
la catégorie abélienne des groupes abéliens, contravariant en $A$ et covariant
en $B$, et exact à gauche par rapport aux deux arguments (c'est-à-dire un
bifoncteur exact à gauche dans la terminologie de [6]).

\LayerTarget{tohoku:section:1.5}
\noindent\textbf{1.5. Sommes et produits infinis.}
Dans certaines constructions, nous aurons besoin de l'existence et de
certaines propriétés des sommes directes et produits directs infinis. Voici
par ordre de force croissante, les axiomes les plus utilisés.

\LayerTarget{tohoku:section:1.5:AB3}
\noindent\textbf{AB 3)} \textit{Pour toute famille $(A_i)_{i\in I}$ d'objet de
$C$, la somme directe des $A_i$ (cf. N\textsuperscript{o}1) existe.}

\LayerTarget{tohoku:section:1.5:p1}
Cet axiome implique que pour toute famille de sous-trucs $A_i$ d'un $A\in C$,
le sup des $A_i$ existe : il suffit de prendre l'image de la somme directe
$\bigoplus_{i\in I}A_i$ par le morphisme dont les composantes sont les
injections canoniques $A_i\longrightarrow A$. Nous avons vu que la somme
directe d'une famille quelconque de morphismes surjectifs est surjectif
(N\textsuperscript{o}1) ; en fait, on voit même que le foncteur

\[
  (A_i)_{i\in I}\longmapsto\bigoplus_{i\in I}A_i,
\]

\noindent défini sur la ``catégorie produit'' $C^I$, et à valeurs

\TohokuContinuation{129}
